\documentclass[11pt]{article}
\usepackage[margin=1in]{geometry}
\usepackage{hyperref}
\usepackage{enumitem}
\title{Species Distribution Models in the Era of Deep Learning: A Synthesis\\
\large A Literature Review Generated with AI Assistance}
\author{LitReview AI \and A. L. Robles Fernández\\[4pt]
\small LitReview \textbar{} ai.litreview.org\\
\small Identifier: 2608.00001}
\date{2026-08-25}

\begin{document}
\maketitle

\begin{abstract}
Species distribution models (SDMs) have evolved from correlative envelope methods to deep learning architectures capable of exploiting massive biodiversity data streams. This review synthesizes progress across four axes: (1) the transition from GLM/GAM-based approaches to neural and ensemble methods, including convolutional and transformer architectures applied to remote sensing covariates; (2) the integration of process-based mechanistic models with data-driven predictions; (3) uncertainty quantification and transferability under climate change scenarios; and (4) emerging evaluation standards for model comparison. We identify key open challenges, including sampling bias correction at continental scales, temporal validity of static models, and the need for interpretable predictions to inform conservation policy.

\noindent\textbf{Keywords:} species distribution models, deep learning, bioclimatic variables, ecological niche.
\end{abstract}

\section*{Placeholder Notice}
This document is a placeholder entry in the LitReview repository. The full
review text will be generated with AI assistance from primary literature and
released after curation. This page establishes the identifier, metadata, and
distribution formats (LaTeX, PDF, HTML) for the entry.

\section*{Scope}
The forthcoming review will synthesize recent literature on the topic above,
covering its principal methods, findings, open challenges, and directions for
future work within the field of Ecology \& Evolution.

\section*{Data and Code Availability}
No new data or code are associated with this placeholder.

\bibliographystyle{plain}
\begin{thebibliography}{9}
\bibitem{litreview2026} LitReview Consortium (2026). \emph{LitReview: an arXiv-style repository of AI-generated literature reviews for the basic sciences.} ai.litreview.org.
\end{thebibliography}

\end{document}
